\documentclass{amsart}
% For dealing with references we use the comment environment
\usepackage{verbatim}
\newenvironment{reference}{\comment}{\endcomment}
%\newenvironment{reference}{}{}
\newenvironment{slogan}{\comment}{\endcomment}
\newenvironment{history}{\comment}{\endcomment}

% For commutative diagrams we use Xy-pic
\usepackage[all]{xy}

% We use 2cell for 2-commutative diagrams.
\xyoption{2cell}
\UseAllTwocells

% We use multicol for the list of chapters between chapters
\usepackage{multicol}

% This is generally recommended for better output
\usepackage{lmodern}
\usepackage[T1]{fontenc}

% For cross-file-references

% Package for hypertext links:
\usepackage{hyperref}

% For any local file, say "hello.tex" you want to link to please

% Theorem environments.
%
\theoremstyle{plain}
\newtheorem{theorem}[subsection]{Theorem}
\newtheorem{proposition}[subsection]{Proposition}
\newtheorem{lemma}[subsection]{Lemma}

\theoremstyle{definition}
\newtheorem{definition}[subsection]{Definition}
\newtheorem{example}[subsection]{Example}
\newtheorem{exercise}[subsection]{Exercise}
\newtheorem{situation}[subsection]{Situation}

\theoremstyle{remark}
\newtheorem{remark}[subsection]{Remark}
\newtheorem{remarks}[subsection]{Remarks}

\numberwithin{equation}{subsection}

% Macros
%
\def\lim{\mathop{\mathrm{lim}}\nolimits}
\def\colim{\mathop{\mathrm{colim}}\nolimits}
\def\Spec{\mathop{\mathrm{Spec}}}
\def\Hom{\mathop{\mathrm{Hom}}\nolimits}
\def\Ext{\mathop{\mathrm{Ext}}\nolimits}
\def\SheafHom{\mathop{\mathcal{H}\!\mathit{om}}\nolimits}
\def\SheafExt{\mathop{\mathcal{E}\!\mathit{xt}}\nolimits}
\def\Sch{\mathit{Sch}}
\def\Mor{\mathop{\mathrm{Mor}}\nolimits}
\def\Ob{\mathop{\mathrm{Ob}}\nolimits}
\def\Sh{\mathop{\mathit{Sh}}\nolimits}
\def\NL{\mathop{N\!L}\nolimits}
\def\CH{\mathop{\mathrm{CH}}\nolimits}
\def\proetale{{pro\text{-}\acute{e}tale}}
\def\etale{{\acute{e}tale}}
\def\QCoh{\mathit{QCoh}}
\def\Ker{\mathop{\mathrm{Ker}}}
\def\Im{\mathop{\mathrm{Im}}}
\def\Coker{\mathop{\mathrm{Coker}}}
\def\Coim{\mathop{\mathrm{Coim}}}

% Boxtimes
%
\DeclareMathSymbol{\boxtimes}{\mathbin}{AMSa}{"02}

%
% Macros for moduli stacks/spaces
%
\def\QCohstack{\mathcal{QC}\!\mathit{oh}}
\def\Cohstack{\mathcal{C}\!\mathit{oh}}
\def\Spacesstack{\mathcal{S}\!\mathit{paces}}
\def\Quotfunctor{\mathrm{Quot}}
\def\Hilbfunctor{\mathrm{Hilb}}
\def\Curvesstack{\mathcal{C}\!\mathit{urves}}
\def\Polarizedstack{\mathcal{P}\!\mathit{olarized}}
\def\Complexesstack{\mathcal{C}\!\mathit{omplexes}}
% \Pic is the operator that assigns to X its picard group, usage \Pic(X)
% \Picardstack_{X/B} denotes the Picard stack of X over B
% \Picardfunctor_{X/B} denotes the Picard functor of X over B
\def\Pic{\mathop{\mathrm{Pic}}\nolimits}
\def\Picardstack{\mathcal{P}\!\mathit{ic}}
\def\Picardfunctor{\mathrm{Pic}}
\def\Deformationcategory{\mathcal{D}\!\mathit{ef}}

\setlength{\emergencystretch}{3em}
\begin{document}
\title[Stacks Project, Tag 0A36]{637 --- Cech vanishing from surjectivity on unions of intersections}
\maketitle
\noindent Copyright (C) 2005--2025 Johan de Jong. Stacks Project authors. Source: \href{https://stacks.math.columbia.edu/tag/0A36}{Tag 0A36}.

\noindent Editorial difficulty note (not author text). Difficulty H2: The proof builds an auxiliary topological space of nonempty index subsets, classifies its open coverings, and transfers the original sheaf to a flasque sheaf with the same Cech complex. This specialized geometric-combinatorial construction is beyond ordinary qualification cohomology..

\noindent Editorial provenance note (not author text). History: excerpt from revision \texttt{a0444\allowbreak 6e57e\allowbreak c1fbc\allowbreak 252a8\allowbreak 71afc\allowbreak ec775\allowbreak 2fb28\allowbreak 07b14}, cohomology.tex, lines 1998--2083. Reference presentation was adapted; original-author context and core support blocks were retained or added. The mathematical wording is unchanged. Licensed under GNU FDL 1.2 or later; no invariant sections or cover texts. See ../licenses/COPYING. Context blocks reproduce author definitions and supporting results; they are not additional counted problems.

\begin{lemma}
\label{lemma-vanishing-ravi}
Let $X$ be a topological space. Let $\mathcal{F}$ be an abelian sheaf
on $X$. Let $\mathcal{U} : U = \bigcup_{i \in I} U_i$ be an
open covering. Assume the restriction mappings
$\mathcal{F}(U) \to \mathcal{F}(U')$ are surjective
for $U'$ an arbitrary union of opens of the form $U_{i_0 \ldots i_p}$.
Then $\check{H}^p(\mathcal{U}, \mathcal{F})$
vanishes for $p > 0$.
\end{lemma}

\begin{proof}
Let $Y$ be the set of nonempty subsets of $I$. We will use the letters
$A, B, C, \ldots$ to denote elements of $Y$, i.e., nonempty subsets of $I$.
For a finite nonempty subset $J \subset I$ let
$$
V_J = \{A \in Y \mid J \subset A\}
$$
This means that $V_{\{i\}} = \{A \in Y \mid i \in A\}$ and
$V_J = \bigcap_{j \in J} V_{\{j\}}$.
Then $V_J \subset V_K$ if and only if $J \supset K$.
There is a unique topology on $Y$ such that the collection of
subsets $V_J$ is a basis for the topology on $Y$. Any open is of the form
$$
V = \bigcup\nolimits_{t \in T} V_{J_t}
$$
for some family of finite subsets $J_t$. If $J_t \subset J_{t'}$
then we may remove $J_{t'}$ from the family without changing $V$.
Thus we may assume there are no inclusions among the $J_t$.
In this case the minimal elements of $V$ are the sets $A = J_t$.
Hence we can read off the family $(J_t)_{t \in T}$ from the open $V$.

\medskip\noindent
We can completely understand open coverings in $Y$. First, because
the elements $A \in Y$ are nonempty subsets of $I$ we have
$$
Y = \bigcup\nolimits_{i \in I} V_{\{i\}}
$$
To understand other coverings, let $V$ be as above and let $V_s \subset Y$
be an open corresponding to the family $(J_{s, t})_{t \in T_s}$. Then
$$
V = \bigcup\nolimits_{s \in S} V_s
$$
if and only if for each $t \in T$ there exists an $s \in S$ and
$t_s \in T_s$ such that $J_t = J_{s, t_s}$. Namely, as the family
$(J_t)_{t \in T}$ is minimal, the minimal element $A = J_t$
has to be in $V_s$ for some $s$, hence $A \in V_{J_{t_s}}$ for some
$t_s \in T_s$. But since $A$ is also minimal in $V_s$ we conclude
that $J_{t_s} = J_t$.

\medskip\noindent
Next we map the set of opens of $Y$ to opens of $X$. Namely, we send
$Y$ to $U$, we use the rule
$$
V_J \mapsto U_J = \bigcap\nolimits_{i \in J} U_i
$$
on the opens $V_J$, and we extend it to arbitrary opens $V$ by the rule
$$
V = \bigcup\nolimits_{t \in T} V_{J_t}
\mapsto
\bigcup\nolimits_{t \in T} U_{J_t}
$$
The classification of open coverings of $Y$ given above shows that
this rule transforms open coverings into open coverings. Thus we obtain
an abelian sheaf $\mathcal{G}$ on $Y$ by setting
$\mathcal{G}(Y) = \mathcal{F}(U)$ and for
$V = \bigcup\nolimits_{t \in T} V_{J_t}$ setting
$$
\mathcal{G}(V) = \mathcal{F}\left(\bigcup\nolimits_{t \in T} U_{J_t}\right)
$$
and using the restriction maps of $\mathcal{F}$.

\medskip\noindent
With these preliminaries out of the way we can prove our lemma as follows.
We have an open covering
$\mathcal{V} : Y = \bigcup_{i \in I} V_{\{i\}}$ of $Y$.
By construction we have an equality
$$
\check{C}^\bullet(\mathcal{V}, \mathcal{G}) =
\check{C}^\bullet(\mathcal{U}, \mathcal{F})
$$
of {\v C}ech complexes. Since the sheaf $\mathcal{G}$ is flasque on $Y$
(by our assumption on $\mathcal{F}$ in the statement of the lemma)
the vanishing follows from
Lemma \href{https://stacks.math.columbia.edu/tag/09SZ}{lemma flasque acyclic cech (Tag 09SZ)}.
\end{proof}
\end{document}
